%% kapa: exam=1 round=29 year=2018 period=1 form=A subject=civil_law subject_ko=민법 seq=19 number=19
%% source: Q-Net 감정평가사 기출문제 (cst00309) artlSeq=5202808 file=extracted/1008/5202808/2018년 제29회 감정평가사 1차시험 기출문제/2018년 제29회 감정평가사 1차시험 기출문제/1교시 A형.pdf
%% conversion: kordoc 4.24.51 -> tools/gen1.py (auto); review=proofread-sonnet; flags=none
\begin{question}{19}
법률행위의 조건과 기한에 관한 설명으로 옳은 것은?
\begin{choices}
  \item 조건은 법률행위의 효력의 발생 또는 소멸을 장래에 생기는 것이 확실한 사실에 의존하게 하는 법률행위의 부관이다.
  \item 법률행위 당시에 곧바로 효력을 발생하게 할 필요가 있는 입양에는 시기를 붙이지 못한다.
  \item 단독행위의 경우 상대방이 동의한 경우에도 조건을 붙일 수 없다.
  \item 정지조건있는 법률행위에서 당사자는 조건성취의 효력을 그 성취전에 소급하게 할 수 없다.
  \item 종기있는 법률행위는 기한이 도래한 때로부터 그 효력이 생긴다.
\end{choices}
\end{question}
\answer{2}
